\section{链同伦映射的构造}
\label{sec:construction-chain-homotopy-map}

虽然上述同调同构只依赖于端点哈密顿量 $H_\alpha,H_\beta$，相应的链映射却不仅依赖
连接 $H_\alpha$ 与 $H_\beta$ 的同伦
$\mathcal H=\{H(\eta)\}_{0\leq\eta\leq1}$，还依赖同伦
$\mathfrak j=\{J(\eta)\}_{0\leq\eta\leq1}$。固定非退化哈密顿量
$H_\alpha,H_\beta$ 以及连接它们的同伦 $\mathcal H$，再固定相容殆复结构的同伦
$\mathfrak j=\{J(\eta)\}_{0\leq\eta\leq1}$ 和截断函数
$\rho:\mathbb R\to[0,1]$。

回忆，在前几节定义模空间
$\mathcal M(H,J;[z^-,w^-],[z^+,w^+])$ 与\allowbreak
$\mathcal M((\mathcal H,\mathfrak j;\rho);[z^-,w^-],[z^+,w^+])$ 时，施加了同伦条件
\begin{equation}
 [z^+,w^+]=[z^+,w^-\#u];\qquad
 [u]=C\quad\text{于 }\pi_2(z^-,z^+).
\label{eq:19-chain-homotopy-class-condition}
\end{equation}
条件~\eqref{eq:19-chain-homotopy-class-condition} 的一个推论是
\[
 [z^+,w^+]=[z^+,w^-\#u]\quad\text{于 }\Gamma,
\]
但后者比前者弱。换言之，可能存在不止一对互异元素
$C_1,C_2\in\pi_2(z^-,z^+)$，使得
\[
 \mu(z^-,z^+;C_1)=\mu(z^-,z^+;C_2),
 \qquad\omega(C_1)=\omega(C_2).
\]

给定同伦的同伦 $(\overline{\mathcal H},\overline{\mathfrak j})$，其中
$\overline{\mathfrak j}=\{\mathfrak j_\kappa\}$、
$\overline{\mathcal H}=\{\mathcal H_\kappa\}$，还定义由
$\mathcal H_\kappa^\rho$ 组成的拉长
\[
 \overline{\mathcal H}^{\rho}
 =\{\mathcal H_\kappa^\rho\}_{0\leq\kappa\leq1}.
\]
对每个 $0\leq\kappa\leq1$，考虑 (18.5.44) 的参数化版本
\begin{equation}
 \left\{
 \begin{aligned}
 &\frac{\partial u}{\partial\tau}
 +J_\kappa^\rho\left(\frac{\partial u}{\partial t}
 -X_{H_\kappa^\rho}(u)\right)=0,\\
 &u(-\infty)=[z^-,w^-],\qquad u(\infty)=[z^+,w^+],
 \end{aligned}
 \right.
\label{eq:19-parameterized-continuation-equation}
\end{equation}
并定义参数化模空间
\[
\begin{aligned}
 &\mathcal M^{\mathrm{para}}
 ((\overline{\mathcal H},\overline{\mathfrak j};\rho);
 [z^-,w^-],[z^+,w^+])\\
 &\quad=\bigcup_{\kappa\in[0,1]}\{\kappa\}\times
 \mathcal M((\mathcal H_\kappa,\mathfrak j_\kappa;\rho);
 [z^-,w^-],[z^+,w^+]).
\end{aligned}
\]
这里要强调，在考察参数化模空间时，渐近条件不依赖参数
$\kappa\in[0,1]$，而是固定的。

考虑这样的对子 $[z^-,w^-],[z^+,w^+]$：相伴参数化模空间满足
\[
 \dim\mathcal M^{\mathrm{para}}
 \bigl((\overline{\mathcal H},\overline{\mathfrak j};\rho);
 [z^-,w^-],[z^+,w^+]\bigr)=0.
\]
借助其基数给出的矩阵元，定义链同伦映射
\footnote{译者注：原书这里把用于计数的空间写成未参数化的
$\mathcal M((\mathcal H_\kappa,\mathfrak j_\kappa;\rho);\cdot,\cdot)$，其中
$\kappa$ 未量化；但前文刚定义参数化模空间，后文又明确由参数空间贡献附加的一维，
故这里依上下文改为 $\mathcal M^{\mathrm{para}}$。}
\[
 \Upsilon_{\overline{\mathcal H}^{\rho}}:
 CF_*(H_\alpha)\longrightarrow CF_*(H_\beta).
\]
（这里仍像此前一样，在半正性假设下使用基于维数计数的横截性--紧性论证。）
由于参数空间贡献一个附加维数，得到
\[
 \mu_{H_\alpha}([z^-,w^-])-\mu_{H_\beta}([z^+,w^+])=-1,
\]
故映射 $\Upsilon_{\overline{\mathcal H}^{\rho}}$ 的次数为 $1$。

再次利用\personGromov{}--\personFloer{}紧性，可得到
$\overline{\mathcal M}^{\mathrm{para}}$ 的边界结构定理。

\begin{proposition}
\label{prop:boundary-parameterized-chain-homotopy-moduli}
假设 $(\mathcal H_0,\mathfrak j_0)$ 与 $(\mathcal H_1,\mathfrak j_1)$ 在定义 19.3.2
的意义下\personFloer{}正则，且 $(\overline{\mathcal H},\overline{\mathfrak j})$ 在相应的
参数化意义下\personFloer{}正则。设
$[z^-,w^-]\in\operatorname{Crit}\mathcal A_{H_\alpha}$ 与
$[z^+,w^+]\in\operatorname{Crit}\mathcal A_{H_\beta}$ 满足
\[
 \mu_{H_\alpha}([z^-,w^-])-\mu_{H_\beta}([z^+,w^+])=0,
\]
从而相伴的参数化模空间维数为 $1$。则对某个整数 $\ell\geq0$，存在有限多个点
$0<\kappa_1,\ldots,\kappa_\ell<1$，并且
\footnote{译者注：源边界式的固定端因子即
$\{0\}\times\mathcal M(H_\alpha,J_\alpha;\cdot)$ 与
$\mathcal M(H_\beta,J_\beta;\cdot)\times\{1\}$，与译文一致；中间段因子源印
覆盖整个 $\kappa$ 区间的 $\mathcal M^{\mathrm{para}}((\overline{\mathcal H},
\overline{\mathfrak j};\rho);\cdot)$。译文按参数化模空间在 $\kappa=0,1$ 端
退化为端点纤维的理解，把中间段改写为
$\mathcal M((\mathcal H_0,\mathfrak j_0;\rho);\cdot)$ 与
$\mathcal M((\mathcal H_1,\mathfrak j_1;\rho);\cdot)$，系解释性重写；源书字面
则保留 $\mathcal M^{\mathrm{para}}$ 因子。末项又把 $\{\kappa_j\}$ 与
$\mathcal H_{\kappa_i}$ 混用，译文统一为 $i$。}
\[
\begin{aligned}
 &\partial\overline{\mathcal M}^{\mathrm{para}}
 ((\overline{\mathcal H},\overline{\mathfrak j};\rho);
 [z^-,w^-],[z^+,w^+])\\
 &=\bigcup_{[z'^-,w'^-]}
 \bigl(\{0\}\times\mathcal M(H_\alpha,J_\alpha;[z^-,w^-],[z'^-,w'^-])\bigr)\\
 &\qquad\times
 \mathcal M((\mathcal H_0,\mathfrak j_0;\rho);
 [z'^-,w'^-],[z^+,w^+])\\
 &\quad\cup\bigcup_{[z'^+,w'^+]}
 \mathcal M((\mathcal H_1,\mathfrak j_1;\rho);
 [z^-,w^-],[z'^+,w'^+])\\
 &\qquad\times
 \bigl(\{1\}\times\mathcal M(H_\beta,J_\beta;[z'^+,w'^+],[z^+,w^+])\bigr)\\
 &\quad\cup\bigcup_{i=1}^{\ell}\{\kappa_i\}\times
 \mathcal M((\mathcal H_{\kappa_i},\mathfrak j_{\kappa_i};\rho);
 [z^-,w^-],[z^+,w^+]).
\end{aligned}
\]
\end{proposition}

于是该命题立即给出如下恒等式（差一个符号）：
\begin{equation}
 h_{(\mathfrak j_1,\mathcal H_1;\rho_1)}
 -h_{(\mathfrak j_0,\mathcal H_0;\rho_0)}
 =\partial_{(J^1,H^1)}\circ\Upsilon_{\overline{\mathcal H}^{\rho}}
 +\Upsilon_{\overline{\mathcal H}^{\rho}}\circ\partial_{(J^0,H^0)}.
\label{eq:19-chain-homotopy-identity}
\end{equation}
映射 $\Upsilon_{\overline{\mathcal H}^{\rho}}$ 仍依赖于同伦
$\overline{\mathfrak j}$ 与 $\rho$ 的选择，故也记作
\[
 \Upsilon_{\overline{\mathcal H}^{\rho}}
 =\Upsilon_{(\overline{\mathcal H},\overline{\mathfrak j};\rho)}.
\]
式~\eqref{eq:19-chain-homotopy-identity} 特别证明了：连接相同端点的不同同伦
$(\mathfrak j_0,\mathcal H_0;\rho)$ 与 $(\mathfrak j_1,\mathcal H_1;\rho)$ 所给出的
两个链映射链同伦。因此，链映射~\eqref{eq:19-floer-homology-continuation-map} \emph{在同调上}
诱导的映射与同伦 $(\mathcal H,\mathfrak j)$ 或 $\rho$ 的选择无关。

\begin{remark}
\label{rem:vary-elongation-in-chain-homotopy}
上面的讨论没有让拉长函数 $\rho$ 随 $\kappa$ 变化，这只是为了简化记号。最自然的
做法其实是让拉长函数依赖于 $\kappa$。
\end{remark}
